\section{Results}\label{results}

\subsection{5.1 Static tail forecasts fail during the 2023 structural
break}\label{static-tail-forecasts-fail-during-the-2023-structural-break}

In the original state-commodity development test, the fixed training q90
was exceeded by approximately 47\% of observations, rather than the
intended 10\%. Static empirical q90 coverage fell to 52.7\%, and the
static LightGBM coverage fell below 47\%. The problem was not primarily
cross-sectional ranking; it was a common upward shift beyond the range
represented in training.

Delayed adaptation materially reduced the loss. The stress-gated EWMA
model achieved pinball loss 4.230 and coverage 71.7\%, compared with
6.836 and 52.7\% for the static empirical q90. The loss reduction was
38.1\%. Rolling conformal calibration also improved performance, but
coverage remained below nominal. These results motivated the later
expert-aggregation design.

\textbf{Table 3. Original locked structural-break test, January
2023-January 2024}

\begin{longtable}[]{@{}
  >{\raggedright\arraybackslash}p{(\columnwidth - 4\tabcolsep) * \real{0.3243}}
  >{\raggedright\arraybackslash}p{(\columnwidth - 4\tabcolsep) * \real{0.3243}}
  >{\raggedright\arraybackslash}p{(\columnwidth - 4\tabcolsep) * \real{0.3243}}@{}}
\toprule\noalign{}
\begin{minipage}[b]{\linewidth}\raggedright
\textbf{Model}
\end{minipage} & \begin{minipage}[b]{\linewidth}\raggedright
\textbf{Pinball loss}
\end{minipage} & \begin{minipage}[b]{\linewidth}\raggedright
\textbf{Coverage}
\end{minipage} \\
\midrule\noalign{}
\endhead
\bottomrule\noalign{}
\endlastfoot
Static empirical q90 & 6.836 & 52.7\% \\
Static quantile LightGBM & 8.338 & 46.6\% \\
Rolling empirical q90, 36 months & 5.634 & 59.8\% \\
Rolling conformal-style calibration & 4.643 & 67.3\% \\
Stress-gated delayed EWMA & 4.230 & 71.7\% \\
\end{longtable}

Notes: The target is the 90th percentile of three-month forward log
food-price inflation. Coverage is the share of outcomes below the
forecast quantile. The 2023 period was locked before the Phase 7
backtest but informed the architecture of subsequent phases.

\subsection{5.2 Extended NBS and external WFP
performance}\label{extended-nbs-and-external-wfp-performance}

\begin{figure}
\centering
\includegraphics{figures/fig2_nbs_model_comparison.pdf}
\caption{Pooled post-2022 q90 pinball loss in the NBS national
continuation.}\label{fig:figure-2}
\end{figure}

\begin{figure}
\centering
\includegraphics{figures/fig3_wfp_model_comparison.pdf}
\caption{Pooled post-2022 q90 pinball loss in the external WFP market
panel.}\label{fig:figure-3}
\end{figure}

\textbf{Table 4. Post-2022 extended and external model comparison}

\begin{longtable}[]{@{}
  >{\raggedright\arraybackslash}p{(\columnwidth - 6\tabcolsep) * \real{0.2368}}
  >{\raggedright\arraybackslash}p{(\columnwidth - 6\tabcolsep) * \real{0.2632}}
  >{\raggedright\arraybackslash}p{(\columnwidth - 6\tabcolsep) * \real{0.2368}}
  >{\raggedright\arraybackslash}p{(\columnwidth - 6\tabcolsep) * \real{0.2368}}@{}}
\toprule\noalign{}
\begin{minipage}[b]{\linewidth}\raggedright
\textbf{Dataset}
\end{minipage} & \begin{minipage}[b]{\linewidth}\raggedright
\textbf{Model}
\end{minipage} & \begin{minipage}[b]{\linewidth}\raggedright
\textbf{Pinball loss}
\end{minipage} & \begin{minipage}[b]{\linewidth}\raggedright
\textbf{Coverage}
\end{minipage} \\
\midrule\noalign{}
\endhead
\bottomrule\noalign{}
\endlastfoot
NBS continuation & Static empirical q90 & 5.582 & 70.2\% \\
NBS continuation & Fixed-persistence gate & 4.218 & 85.7\% \\
NBS continuation & Global conformal (12m) & 5.052 & 72.6\% \\
NBS continuation & Pure hierarchical conformal & 4.481 & 73.4\% \\
NBS continuation & Ungated expert mixture & 4.459 & 77.9\% \\
NBS continuation & Meta-controller & 4.206 & 82.9\% \\
WFP external panel & Static empirical q90 & 6.679 & 85.0\% \\
WFP external panel & Fixed-persistence gate & 6.494 & 86.6\% \\
WFP external panel & Global conformal (12m) & 6.340 & 89.2\% \\
WFP external panel & Pure hierarchical conformal & 6.397 & 85.5\% \\
WFP external panel & Ungated expert mixture & 6.089 & 88.4\% \\
WFP external panel & Meta-controller & 6.111 & 88.3\% \\
\end{longtable}

Notes: The primary controller is the stress-override inertial expert
mixture selected jointly on the 2022 validation objective. The ungated
mixture has slightly lower WFP loss but is materially weaker in the NBS
continuation.

In the NBS continuation, the meta-controller reduces pooled pinball loss
from 5.582 for the static q90 to 4.206 and closes 74.5\% of the
equal-month gap between the static forecast and the monthly hindsight
oracle. Its coverage is 82.9\%, substantially improved but still below
the nominal 90\%. The fixed-persistence gate is nearly as accurate in
pooled loss (4.218) and has higher coverage (85.7\%).

In the external WFP panel, the controller reduces loss from 6.679 to
6.111 and closes 67.2\% of the static-to-oracle gap. Coverage increases
from 85.0\% to 88.3\%. Global conformal calibration attains the best
individual-expert coverage (89.2\%) and loss 6.340. The ungated mixture
has slightly lower pooled WFP loss than the primary controller, but the
stress-override specification is the one selected jointly across both
sources using the frozen validation objective.

\subsection{5.3 Bootstrap comparisons and statistical
interpretation}\label{bootstrap-comparisons-and-statistical-interpretation}

\textbf{Table 5. Equal-month bootstrap comparisons for the
meta-controller}

\begin{longtable}[]{@{}
  >{\raggedright\arraybackslash}p{(\columnwidth - 8\tabcolsep) * \real{0.1948}}
  >{\raggedright\arraybackslash}p{(\columnwidth - 8\tabcolsep) * \real{0.1948}}
  >{\raggedright\arraybackslash}p{(\columnwidth - 8\tabcolsep) * \real{0.1948}}
  >{\raggedright\arraybackslash}p{(\columnwidth - 8\tabcolsep) * \real{0.1948}}
  >{\raggedright\arraybackslash}p{(\columnwidth - 8\tabcolsep) * \real{0.1948}}@{}}
\toprule\noalign{}
\begin{minipage}[b]{\linewidth}\raggedright
\textbf{Dataset}
\end{minipage} & \begin{minipage}[b]{\linewidth}\raggedright
\textbf{Baseline}
\end{minipage} & \begin{minipage}[b]{\linewidth}\raggedright
\textbf{Mean loss difference}
\end{minipage} & \begin{minipage}[b]{\linewidth}\raggedright
\textbf{95\% interval}
\end{minipage} & \begin{minipage}[b]{\linewidth}\raggedright
\textbf{p-value}
\end{minipage} \\
\midrule\noalign{}
\endhead
\bottomrule\noalign{}
\endlastfoot
NBS & Static q90 & -1.366 & {[}-3.000, -0.027{]} & 0.045 \\
NBS & Fixed gate & -0.012 & {[}-0.131, 0.133{]} & 0.810 \\
WFP & Static q90 & -0.633 & {[}-0.957, -0.327{]} & 0.000 \\
WFP & Global conformal & -0.166 & {[}-0.474, 0.074{]} & 0.201 \\
\end{longtable}

Notes: Negative differences favour the meta-controller. The bootstrap
resamples forecast-origin months and is descriptive because the
controller architecture was motivated after earlier examination of the
same post-2022 period.

The bootstrap rejects equal performance against the static forecast in
both datasets: the NBS 95\% interval is {[}-3.000, -0.027{]} with
\(p\approx 0.045\), while the WFP interval is {[}-0.957, -0.327{]} with
p\textless0.001. In contrast, the interval against the best expert
includes zero in both datasets. The controller is therefore meaningfully
better than remaining static, but the data do not establish confirmatory
superiority over the strongest individual adaptive rule.

\subsection{5.4 Regime-specific performance and recovery
failure}\label{regime-specific-performance-and-recovery-failure}

\textbf{Table 6. NBS performance across break and recovery periods}

\begin{longtable}[]{@{}
  >{\raggedright\arraybackslash}p{(\columnwidth - 6\tabcolsep) * \real{0.2432}}
  >{\raggedright\arraybackslash}p{(\columnwidth - 6\tabcolsep) * \real{0.2432}}
  >{\raggedright\arraybackslash}p{(\columnwidth - 6\tabcolsep) * \real{0.2432}}
  >{\raggedright\arraybackslash}p{(\columnwidth - 6\tabcolsep) * \real{0.2432}}@{}}
\toprule\noalign{}
\begin{minipage}[b]{\linewidth}\raggedright
\textbf{Period}
\end{minipage} & \begin{minipage}[b]{\linewidth}\raggedright
\textbf{Model}
\end{minipage} & \begin{minipage}[b]{\linewidth}\raggedright
\textbf{Pinball loss}
\end{minipage} & \begin{minipage}[b]{\linewidth}\raggedright
\textbf{Coverage}
\end{minipage} \\
\midrule\noalign{}
\endhead
\bottomrule\noalign{}
\endlastfoot
Jan 2023-Jan 2024 & Static & 9.494 & 48.7\% \\
Jan 2023-Jan 2024 & Fixed gate & 6.000 & 72.4\% \\
Jan 2023-Jan 2024 & Meta-controller & 6.161 & 66.6\% \\
Feb-Dec 2024 & Static & 3.619 & 70.5\% \\
Feb-Dec 2024 & Fixed gate & 3.392 & 97.4\% \\
Feb-Dec 2024 & Meta-controller & 3.259 & 97.4\% \\
2025 & Static & 1.469 & 96.9\% \\
2025 & Fixed gate & 2.466 & 99.1\% \\
2025 & Meta-controller & 2.276 & 99.1\% \\
2026 available origins & Static & 0.922 & 92.9\% \\
2026 available origins & Fixed gate & 0.922 & 92.9\% \\
2026 available origins & Meta-controller & 0.778 & 92.9\% \\
\end{longtable}

Notes: The 2026 extension contains only one usable forecast-origin month
and should not be interpreted as a stable annual comparison.

During January 2023-January 2024, the static NBS loss is 9.494, compared
with 6.000 for the fixed gate and 6.161 for the controller. The
controller is not the best model during the most acute break because its
mixture retains weight on weaker experts. During February-December 2024
it improves to 3.259, modestly below the static loss of 3.619.

The main recovery failure appears in 2025. The static forecast records
loss 1.469, whereas the controller records 2.276 and the fixed gate
2.466. Coverage near 99\% reveals overconservatism rather than improved
risk control. The inherited six-month stress persistence remains active
after the tail level has normalised. This is the clearest reason not to
claim universal superiority.

\subsection{5.5 Dynamic regret}\label{dynamic-regret}

\begin{figure}
\centering
\includegraphics{figures/fig4_nbs_cumulative_regret.pdf}
\caption{Cumulative equal-month regret relative to the monthly hindsight
oracle, NBS continuation.}\label{fig:figure-4}
\end{figure}

\begin{figure}
\centering
\includegraphics{figures/fig5_wfp_cumulative_regret.pdf}
\caption{Cumulative equal-month regret relative to the monthly hindsight
oracle, WFP market panel.}\label{fig:figure-5}
\end{figure}

The cumulative-regret paths show that the controller's principal gain is
avoidance of the static model's large 2023 error accumulation. In NBS,
its equal-month dynamic regret is 0.468 compared with 1.834 for the
static forecast and 0.481 for the fixed gate. In WFP, the corresponding
values are 0.309, 0.942, and 0.475 for global conformal. The controller
smooths expert-specific failures, but it does not track the monthly
oracle closely enough to eliminate regret.

\subsection{5.6 Conditional coverage}\label{conditional-coverage}

Pooled performance conceals subgroup failures. In the NBS continuation,
the controller's worst commodity coverage is 65.2\% for Irish potato. In
WFP, the worst commodity coverage is 80.9\% for tomatoes and the worst
market coverage is 83.6\% for Monday market. Global conformal provides
stronger WFP mean commodity coverage, while the fixed gate provides
stronger NBS mean commodity coverage. These results rule out a claim of
conditional calibration and show why pooled q90 coverage cannot be the
only diagnostic.

\subsection{5.7 Simulation evidence}\label{simulation-evidence}

\begin{figure}
\centering
\includegraphics{figures/fig6_simulation_regret.pdf}
\caption{Mean dynamic regret in delayed-feedback structural-break
simulations.}\label{fig:figure-6}
\end{figure}

\textbf{Table 7. Lowest-loss method by simulated scenario}

\begin{longtable}[]{@{}
  >{\raggedright\arraybackslash}p{(\columnwidth - 6\tabcolsep) * \real{0.2368}}
  >{\raggedright\arraybackslash}p{(\columnwidth - 6\tabcolsep) * \real{0.2632}}
  >{\raggedright\arraybackslash}p{(\columnwidth - 6\tabcolsep) * \real{0.2368}}
  >{\raggedright\arraybackslash}p{(\columnwidth - 6\tabcolsep) * \real{0.2368}}@{}}
\toprule\noalign{}
\begin{minipage}[b]{\linewidth}\raggedright
\textbf{Scenario}
\end{minipage} & \begin{minipage}[b]{\linewidth}\raggedright
\textbf{Lowest-loss method}
\end{minipage} & \begin{minipage}[b]{\linewidth}\raggedright
\textbf{Mean pinball}
\end{minipage} & \begin{minipage}[b]{\linewidth}\raggedright
\textbf{Coverage}
\end{minipage} \\
\midrule\noalign{}
\endhead
\bottomrule\noalign{}
\endlastfoot
Abrupt Permanent & Pure hierarchical conformal & 0.884 & 79.9\% \\
Gradual Break Recovery & Pure hierarchical conformal & 0.832 & 79.9\% \\
Heterogeneous Subgroups & Pure hierarchical conformal & 1.039 &
86.8\% \\
No Break & Static & 0.534 & 90.1\% \\
Temporary Recovery & Pure hierarchical conformal & 0.990 & 82.1\% \\
Volatility Break Recovery & Fixed-persistence gate & 0.759 & 88.6\% \\
\end{longtable}

Notes: Each scenario contains 60 replications, 60 series in six groups,
96 months, Student-t disturbances, and a three-month feedback delay.
Scenario winners are descriptive features of the specified
data-generating processes, not universal optimality statements.

No method dominates. The static quantile is best without a break.
Hierarchical calibration is best under permanent, gradual, temporary,
and subgroup-specific level shifts, although its coverage is often below
90\%. The fixed gate is best under the volatility-break recovery
scenario. The controller is a compromise forecast with lower regret than
static forecasting in most changing environments but is not the
lowest-loss method in any primary simulation scenario.
